\tikzsetfigurename{00_preface-}
%% 00_preface.tex -- Preface (front matter, unnumbered). Author: p1
\chapter{Preface}\label{ch:preface}

Every fast interface, from a PCIe slot to the backplane of an Ethernet switch, moves its data over a \textbf{serial link}: a transmitter, a lossy wire and a receiver that recovers both the bits and the clock. This book teaches you to design one by following a single, concrete chip, the \textbf{\OS{}} transceiver: a \qty{5}{\giga\bit\per\second} NRZ link built in the open-source IHP \sgx{} \qty{130}{\nano\meter} SiGe BiCMOS process~\cite{ihppdk} with open-source tools~\cite{iicosic}.

\paragraph{Who this book is for} Third- and fourth-year electrical engineering undergraduates who will help design the \OS{} chip, and anyone who wants a practical first course in high-speed links. After reading it you should be able to:
\begin{itemize}
  \item predict the eye a channel will produce from its loss and pulse response;
  \item choose and size the TX FFE, CTLE and DFE for a channel, and defend the split;
  \item turn a link budget into specifications for the driver, front end, samplers, PLL and CDR;
  \item model the link in Python, then carry it through schematic, layout and sign-off with open tools.
\end{itemize}

\paragraph{Prerequisites} Basic circuits (KCL, small-signal MOS and BJT models, the differential pair), signals and systems (Fourier, Laplace, convolution) and some Python. Transmission lines, equalization, jitter, PLLs and CDRs are built up from first principles.

\section*{How the book is organized}
The book has five parts, shown in \cref{fig:pre-roadmap}. \textbf{Part~I, Foundations} (Chapters~1--4), gives you the vocabulary of links (pulse response, eye diagram, bit-error rate), the physics of the channel and the tools for a link budget. \textbf{Part~II, Equalization} (Chapters~5--8), shows how the TX FFE, the CTLE and the DFE undo the channel, and how to split the work between them. \textbf{Part~III, Transceiver Circuits} (Chapters~9--10), turns the system model into HBT and CMOS circuits. \textbf{Part~IV, Clocking} (Chapters~11--13), covers jitter, the PLL and clock and data recovery. \textbf{Part~V, Putting It Together} (Chapters~14--16), assembles the full architecture, walks through the open-source design flow and collects the hands-on labs.

\begin{figure}[htbp]
  \centering
  \includegraphics[width=\linewidth]{canva_roadmap}
  \caption{\textbf{Your path through this tutorial.} The five parts are milestones: Foundations, Equalization, Transceiver circuits, Clocking, and Putting it together. Each part builds on the numbers and models of the one before it.}
  \label{fig:pre-roadmap}
\end{figure}

\paragraph{Two reading paths} Everyone should read Part~I first. After that, \cref{fig:pre-paths} suggests two routes. The \textbf{system and modeling path} suits readers who will own the link budget and the serdespy model. The \textbf{circuit path} suits readers who will design transistors and layout. Both paths meet again at the architecture chapter and the labs.

\begin{figure}[htbp]
\centering
\tikzset{prech/.style={circle,draw=#1,fill=white,line width=0.8pt,minimum size=6.4mm,inner sep=0pt,font=\sffamily\footnotesize\bfseries,text=#1},
  prelk/.style={line width=1.1pt,draw=#1}}
\begin{tikzpicture}[x=1cm,y=1cm]
% trunk
\foreach \n [count=\i] in {1,2,3,4}{\node[prech=navy] (t\n) at ({0.95*(\i-1)},0) {\n};}
\foreach \a/\b in {1/2,2/3,3/4}{\draw[prelk=navy] (t\a) -- (t\b);}
% system path (top)
\foreach \n [count=\i] in {5,6,8,11,13}{\node[prech=bblue] (s\n) at ({3.65+0.95*\i},0.75) {\n};}
\draw[prelk=bblue] (t4) -- (s5); \foreach \a/\b in {5/6,6/8,8/11,11/13}{\draw[prelk=bblue] (s\a) -- (s\b);}
% circuit path (bottom)
\foreach \n [count=\i] in {7,9,10,12,15}{\node[prech=golddk] (c\n) at ({3.65+0.95*\i},-0.75) {\n};}
\draw[prelk=golddk] (t4) -- (c7); \foreach \a/\b in {7/9,9/10,10/12,12/15}{\draw[prelk=golddk] (c\a) -- (c\b);}
% merge
\node[prech=navy] (m14) at (10.0,0) {14};
\node[prech=navy] (m16) at (10.95,0) {16};
\draw[prelk=bblue] (s13) -- (m14); \draw[prelk=golddk] (c15) -- (m14); \draw[prelk=navy] (m14) -- (m16);
\node[font=\sffamily\scriptsize,text=navy,anchor=north] at (1.42,-0.42) {Part I: everyone};
\node[font=\sffamily\scriptsize\bfseries,text=bblue,anchor=south] at (6.5,1.1) {system and modeling path};
\node[font=\sffamily\scriptsize\bfseries,text=golddk,anchor=north] at (6.5,-1.1) {circuit path};
\node[font=\sffamily\scriptsize,text=navy,anchor=north,align=center] at (10.48,-0.42) {architecture\\[-1pt]and labs};
\end{tikzpicture}
\caption{\textbf{Two ways to read the book after Part~I.} Circles are chapter numbers. The system path (blue) follows the equalization theory, jitter and the CDR; the circuit path (gold) follows the CTLE circuits, the TX and RX, the PLL and the design flow. Both meet at the architecture (\cref{ch:arch}) and the labs (\cref{ch:labs}).}
\label{fig:pre-paths}
\end{figure}

\section*{How to read the boxes}
Chapters open with what you will learn and close with key takeaways. In between, six kinds of boxes mark the role of each piece of text:

{\renewcommand{\thewexample}{P.1}%
\begin{keyeq}[Bit period]
\[ \Tb = \frac{1}{\fb} = \frac{1}{\qty{5}{\giga\bit\per\second}} = \qty{200}{\pico\second} \]
\end{keyeq}
\begin{example}{A number you will use on every page}
At \qty{5}{\giga\bit\per\second}, $\fN = \fb/2 = \qty{2.5}{\giga\hertz}$; channel loss is quoted there.
\end{example}}
\setcounter{wexample}{0}
\begin{insight}
An \emph{insight} box says what a result means for the \OS{} design.
\end{insight}
\begin{pitfall}
A \emph{pitfall} box warns about a common mistake, such as quoting loss at \qty{5}{\giga\hertz}.
\end{pitfall}
\begin{tryit}[Lab 0: set-up]
A \emph{try-it} box points to a lab script (\cref{sec:lab0}).
\end{tryit}
\begin{specnote}
A \emph{specnote} box states an \OS{} target: \qty{5}{\giga\bit\per\second} NRZ at $\BER \le 10^{-12}$.
\end{specnote}
Longer algebra sits in grey \emph{derivation} boxes.

\section*{How the labs work}
The labs of \cref{ch:labs} reproduce the key plots of the book with \tool{serdespy}~\cite{serdespy}, a Python library for system-level SerDes modeling, plus NumPy and SciPy. The scripts live in \code{tutorial/code/} and share one file, \code{common.py}, that defines the link, the reference channel and the helpers; the figure scripts in \code{tutorial/figsrc/} import the same model, so labs and figures use the same numbers. Run a lab, change a parameter and compare your plot with the text.

\section*{Sources and acknowledgments}
The link-modeling chapters follow the structure of Prof.\ Sam Palermo's ECEN~720 \emph{High-Speed Links} course at Texas A\&M University~\cite{palermo720}. The equalizer and CDR chapters draw on Benjamin Liu's lecture notes~\cite{liuEq,liuCDR}. The labs are built on serdespy by Richard Barrie and colleagues at the University of Toronto~\cite{serdespy}. The chip itself is possible because IHP opened its \sgx{} process design kit~\cite{ihppdk} and the Institute for Integrated Circuits at JKU Linz maintains the IIC-OSIC-TOOLS container and the analog chip template~\cite{iicosic,jkutemplate}.

This is a \textbf{living draft (v0.1)}; numbers will be refined as the design matures. If you find an error, an unclear derivation or a figure that does not explain what it should, please email \href{mailto:bricd.org@gmail.com}{bricd.org@gmail.com} or open an issue or a pull request on the project's GitHub repository, linked from thebricd.org.

\section*{How this book was made}
This book, its figures and lab code, and the accompanying slides were developed with the help of AI tools. The authors reviewed the text, figures and code, but errors may remain; please report any you find as described above, and we will correct them. The topics, structure and many of the results follow the sources acknowledged above and listed in the References.
